% Options for packages loaded elsewhere
% Options for packages loaded elsewhere
\PassOptionsToPackage{unicode}{hyperref}
\PassOptionsToPackage{hyphens}{url}
\PassOptionsToPackage{dvipsnames,svgnames,x11names}{xcolor}
%
\documentclass[
  twocolumn,
  landscape]{report}
\usepackage{xcolor}
\usepackage{amsmath,amssymb}
\setcounter{secnumdepth}{5}
\usepackage{iftex}
\ifPDFTeX
  \usepackage[T1]{fontenc}
  \usepackage[utf8]{inputenc}
  \usepackage{textcomp} % provide euro and other symbols
\else % if luatex or xetex
  \usepackage{unicode-math} % this also loads fontspec
  \defaultfontfeatures{Scale=MatchLowercase}
  \defaultfontfeatures[\rmfamily]{Ligatures=TeX,Scale=1}
\fi
\usepackage{lmodern}
\ifPDFTeX\else
  % xetex/luatex font selection
\fi
% Use upquote if available, for straight quotes in verbatim environments
\IfFileExists{upquote.sty}{\usepackage{upquote}}{}
\IfFileExists{microtype.sty}{% use microtype if available
  \usepackage[]{microtype}
  \UseMicrotypeSet[protrusion]{basicmath} % disable protrusion for tt fonts
}{}
\makeatletter
\@ifundefined{KOMAClassName}{% if non-KOMA class
  \IfFileExists{parskip.sty}{%
    \usepackage{parskip}
  }{% else
    \setlength{\parindent}{0pt}
    \setlength{\parskip}{6pt plus 2pt minus 1pt}}
}{% if KOMA class
  \KOMAoptions{parskip=half}}
\makeatother
% Make \paragraph and \subparagraph free-standing
\makeatletter
\ifx\paragraph\undefined\else
  \let\oldparagraph\paragraph
  \renewcommand{\paragraph}{
    \@ifstar
      \xxxParagraphStar
      \xxxParagraphNoStar
  }
  \newcommand{\xxxParagraphStar}[1]{\oldparagraph*{#1}\mbox{}}
  \newcommand{\xxxParagraphNoStar}[1]{\oldparagraph{#1}\mbox{}}
\fi
\ifx\subparagraph\undefined\else
  \let\oldsubparagraph\subparagraph
  \renewcommand{\subparagraph}{
    \@ifstar
      \xxxSubParagraphStar
      \xxxSubParagraphNoStar
  }
  \newcommand{\xxxSubParagraphStar}[1]{\oldsubparagraph*{#1}\mbox{}}
  \newcommand{\xxxSubParagraphNoStar}[1]{\oldsubparagraph{#1}\mbox{}}
\fi
\makeatother


\usepackage{longtable,booktabs,array}
\usepackage{calc} % for calculating minipage widths
% Correct order of tables after \paragraph or \subparagraph
\usepackage{etoolbox}
\makeatletter
\patchcmd\longtable{\par}{\if@noskipsec\mbox{}\fi\par}{}{}
\makeatother
% Allow footnotes in longtable head/foot
\IfFileExists{footnotehyper.sty}{\usepackage{footnotehyper}}{\usepackage{footnote}}
\makesavenoteenv{longtable}
\usepackage{graphicx}
\makeatletter
\newsavebox\pandoc@box
\newcommand*\pandocbounded[1]{% scales image to fit in text height/width
  \sbox\pandoc@box{#1}%
  \Gscale@div\@tempa{\textheight}{\dimexpr\ht\pandoc@box+\dp\pandoc@box\relax}%
  \Gscale@div\@tempb{\linewidth}{\wd\pandoc@box}%
  \ifdim\@tempb\p@<\@tempa\p@\let\@tempa\@tempb\fi% select the smaller of both
  \ifdim\@tempa\p@<\p@\scalebox{\@tempa}{\usebox\pandoc@box}%
  \else\usebox{\pandoc@box}%
  \fi%
}
% Set default figure placement to htbp
\def\fps@figure{htbp}
\makeatother





\setlength{\emergencystretch}{3em} % prevent overfull lines

\providecommand{\tightlist}{%
  \setlength{\itemsep}{0pt}\setlength{\parskip}{0pt}}



 
\usepackage[]{biblatex}


\usepackage[utf8]{inputenc}
\usepackage[T1]{fontenc}
\usepackage{lmodern}
\usepackage{amsmath, amssymb, amsthm}
\usepackage{geometry}
\geometry{margin=1in}
\usepackage{enumitem}
\usepackage{graphicx}
\graphicspath{ {./images/} }
\usepackage{mathtools}
\usepackage{physics}
\usepackage{xcolor}
\newtheorem{theorem}{Theorem}[section]
\newtheorem{lemma}[theorem]{Lemma}
\newtheorem{proposition}[theorem]{Proposition}
\newtheorem{corollary}[theorem]{Corollary}
\theoremstyle{definition}
\newtheorem{definition}[theorem]{Definition}
\theoremstyle{remark}
\newtheorem{remark}[theorem]{Remark}
\newtheorem{example}[theorem]{Example}
\newtheorem{claim}[theorem]{Claim}
\newcommand{\R}{\mathbb{R}}
\newcommand{\N}{\mathbb{N}}
\newcommand{\Z}{\mathbb{Z}}
\newcommand{\Q}{\mathbb{Q}}
\newcommand{\C}{\mathbb{C}}
\newcommand{\eps}{\varepsilon}
\def\*#1{\mathbf{#1}}
\makeatletter
\@ifpackageloaded{caption}{}{\usepackage{caption}}
\AtBeginDocument{%
\ifdefined\contentsname
  \renewcommand*\contentsname{Table of contents}
\else
  \newcommand\contentsname{Table of contents}
\fi
\ifdefined\listfigurename
  \renewcommand*\listfigurename{List of Figures}
\else
  \newcommand\listfigurename{List of Figures}
\fi
\ifdefined\listtablename
  \renewcommand*\listtablename{List of Tables}
\else
  \newcommand\listtablename{List of Tables}
\fi
\ifdefined\figurename
  \renewcommand*\figurename{Figure}
\else
  \newcommand\figurename{Figure}
\fi
\ifdefined\tablename
  \renewcommand*\tablename{Table}
\else
  \newcommand\tablename{Table}
\fi
}
\@ifpackageloaded{float}{}{\usepackage{float}}
\floatstyle{ruled}
\@ifundefined{c@chapter}{\newfloat{codelisting}{h}{lop}}{\newfloat{codelisting}{h}{lop}[chapter]}
\floatname{codelisting}{Listing}
\newcommand*\listoflistings{\listof{codelisting}{List of Listings}}
\makeatother
\makeatletter
\makeatother
\makeatletter
\@ifpackageloaded{caption}{}{\usepackage{caption}}
\@ifpackageloaded{subcaption}{}{\usepackage{subcaption}}
\makeatother
\usepackage{bookmark}
\IfFileExists{xurl.sty}{\usepackage{xurl}}{} % add URL line breaks if available
\urlstyle{same}
\hypersetup{
  pdftitle={My document},
  colorlinks=true,
  linkcolor={blue},
  filecolor={Maroon},
  citecolor={Blue},
  urlcolor={Blue},
  pdfcreator={LaTeX via pandoc}}


\title{My document}
\author{}
\date{}
\begin{document}
\maketitle

\renewcommand*\contentsname{Table of contents}
{
\hypersetup{linkcolor=}
\setcounter{tocdepth}{2}
\tableofcontents
}
\listoffigures
\listoftables

These notes are primarily based on Section 4.5 in \cite{Anderson2010}.

This handout is based on Section 2.5, 2.6, and 4.5 in
\cite{Anderson2010}, and on Sections 13 and 14 in \cite{Kemp2022}. We
also refer to \cite{Dimitriu2018}. Throughout this handout, \(Z\) has
referred to a constant whose value makes the associated term a
probability density function or measure. Furthermore, we considered
random elements on the probability space \((\Omega, \mathcal{F}, P)\).
We also refered to \(\mathcal{N}(\mu, \sigma)\) as the Gaussian
distribution on \(\mathbb R\) with mean \(\mu \in \mathbb R\) and
variance \(\sigma > 0\). And the entry in the \(i\)th row and \(j\)th
column of a matrix \(A\) was denoted \(A_{ij}\).

\section{Notation}

Let \(\R_+ = (0,\infty)\).

Let \(O(N) = \{O \in \mathbb R^{N \times N}: OO^T = I_N \}\). We define
\[
O(N)^+ = \{ O \in O(N): O_{1,i} > 0 \text{ for all } i=1, \cdots, N \}.
\] Similarly, we have \[
O(N)_+ = \{ O \in O(N): \text{first nonzero entry in each column of $O$ is positive} \}.
\] Let \(H_N^{(1)}\) denote the set of all \(N \times N\) real symmetric
matrices.

Let \(T_N^{(1)}\) denote the tridiagonal matrices with real entries of
size \(N \times N\). Let \(\mathcal{T}_N\) denote the subset of
tridiagonal matrices satisfying positivity on the off-diagonal. That is,
if \((a_N, \cdots, a_1)\) and \((b_{N-1}, \cdots, b_1)\) characterize
\(T \in \mathcal{T}_N\), then \(b_i > 0\) for all \(i\). We denote \(k\)
dimensional Lebesgue measure on \(\R^{k}\) by \(\mathcal{L}_{k}\). But
if the dummy variable of integration is, say, \(u\), then we would also
write \(\int du\) for integration with respect to Lebesgue measure, and
the dimension can be read off from the context.

We define \(\Delta_N^c = \{x_1 > x_2 > \cdots > x_N\} \subset \R^N\). We
let \[S^{N-1} = \{ x \in \R^N: \lVert x \rVert = 1 \},\] and
\[S_+^{N-1} = \{ x \in S^{N-1}: x_i > 0 \text{ for } i = 1, \cdots, N\}.\]

If we write \(D \in \Delta_N^c\), this means that \(D\) is an
\(N \times N\) diagonal matrix with diagonal entries forming an element
of \(\Delta_N^c\). In this way, we identify \(d \in \Delta_N^c\) with a
diagonal matrix \(D\).

\section{The $\beta$ Ensembles}

The \(\beta\) ensembles allow for presenting different matrix models all
at once. The case \(\beta =1\) corresponds to the Gaussian Orthogonal
Ensemble. The case \(\beta = 2\) corresponds to the Gaussian Unitary
Ensemble. Interestingly, there is a matrix model, the so called
tridiagonal matrix model, which corresponds to any \(\beta\). Is this
useful for physics models where \(\beta \to \infty\)? The author of this
handout is unsure.

\begin{definition}[Statistics of $N$ particle gas, $\beta$-ensemble, $P_{V,\beta}^N$]
Let $\lambda = (\lambda_1, \cdots, \lambda_N)$ be $N$ real-valued random variables. 
We define $P_{V,\beta}^N$ as the probability measure on $\mathbb R^N$ whose
density is proportional to
$$
| \Delta(\lambda) |^\beta \exp \left( -N \sum_{i=1}^N V(\lambda_i) \right) = 
\exp \left( \beta \sum_{1 \le i < j \le N} \log |\lambda_j - \lambda_i| -N \sum_{i=1}^N V(\lambda_i) \right),
$$
where the above equality holds when the $\lambda_i$ are all distinct. 
In the above, the parameter $\beta$ is a real number greater than 0, 
the function $V: \mathbb R \to \mathbb R$ satisfies continuity and logarithmic growth
at infinity:
for some $\beta' \ge \beta$ where $\beta' > 1$, we have 
$$
\liminf_{|x| \to \infty} \frac{V(x)}{\beta' \log|x|} > 1. 
$$
We will discuss how integrability of the 
density function depends on this condition.
\end{definition}


\printbibliography



\end{document}
